% GENERATED FILE. Edit canonical lessons, then run npm run generate:books.
% canonical-source-hash: fnv1a64:9a50d3941b51f46e
% canonical-lessons: IT-C384-norma, IT-C384-regolamento, IT-C384-obbligatorio, IT-C384-facoltativo, IT-C384-obbligo

\chapter{Le regole}
\label{ch:it-le-regole}
\hlchaptermodality{384}

\begin{chapteropening}
\textbf{By the end of this chapter:} \emph{I can talk about rules: a norm, the regulations, what is compulsory or optional, and an obligation.}

Complete the last lesson of chapter 384: I can talk about rules: a norm, the regulations, what is compulsory or optional, and an obligation.
\end{chapteropening}

\section[la norma]{la norma — the rule, the standard, the norm}

\label{lesson:IT-C384-norma}



\begin{quote}
Before the new one: say the Italian for the assessment, then the Italian for the score.
\end{quote}

\subsection*{You'll want to know: la norma}
\textbf{la norma} — \textquotedblleft{}the rule, the standard, the norm\textquotedblright{}.

\begin{quote}
Le norme di sicurezza sono cambiate.
\end{quote}

\begin{quote}
\emph{The safety rules have changed.}
\end{quote}

\emph{Di norma} — as a rule, usually. \emph{A norma di legge} — in line with the law.

\begin{cousinweb}
\emph{norma} is a learned borrowing from Latin \emph{nōrma} \textquotedblleft{}a carpenter's square; a rule, a pattern\textquotedblright{} (Lewis \& Short): the tool that keeps a corner straight.
\end{cousinweb}

\begin{grammarlens}[title={What you've built}]
One more word for telling, explaining and describing at B1.
\end{grammarlens}

\subsection*{Your turn}
\begin{itemize}
\raggedright
  \item \emph{Say it:} \emph{la norma}
  \item \emph{Say it:} \emph{la norma}, once more
  \item \emph{Say it:} a sentence of your own with \emph{la norma}
  \item \emph{Recall:} say the Italian for the assessment, then the Italian for the score, then say \emph{la norma} again
\end{itemize}

\begin{tcolorbox}[breakable,colback=teal!4,colframe=teal!35!black,title={Before you move on}]
What does la norma mean? (\textquotedblleft{}The rule, the standard, the norm\textquotedblright{}.) Say it once more.
\end{tcolorbox}
\section[il regolamento]{il regolamento — the regulations, the rules (of a place)}

\label{lesson:IT-C384-regolamento}



\begin{quote}
Before the new one: say the Italian for the score, then the Italian for the rule.
\end{quote}

\subsection*{You'll want to know: il regolamento}
\textbf{il regolamento} — \textquotedblleft{}the regulations, the rules (of a place)\textquotedblright{}.

\begin{quote}
Secondo il regolamento, in piscina bisogna fare la doccia prima di entrare.
\end{quote}

\begin{quote}
\emph{According to the rules, you must shower before going into the pool.}
\end{quote}

The written rules of a school, a club or a building.

\begin{cousinweb}
\emph{regolamento} is built on \emph{regolare} \textquotedblleft{}to regulate\textquotedblright{}, from Late Latin \emph{rēgulāre}, built on \emph{rēgula} \textquotedblleft{}a ruler, a rule\textquotedblright{} (Lewis \& Short), from \emph{regere} \textquotedblleft{}to guide\textquotedblright{}.
\end{cousinweb}

\begin{grammarlens}[title={What you've built}]
One more word for telling, explaining and describing at B1.
\end{grammarlens}

\subsection*{Your turn}
\begin{itemize}
\raggedright
  \item \emph{Say it:} \emph{il regolamento}
  \item \emph{Say it:} \emph{il regolamento}, once more
  \item \emph{Say it:} a sentence of your own with \emph{il regolamento}
  \item \emph{Recall:} say the Italian for the score, then the Italian for the rule, then say \emph{il regolamento} again
\end{itemize}

\begin{tcolorbox}[breakable,colback=teal!4,colframe=teal!35!black,title={Before you move on}]
What does il regolamento mean? (\textquotedblleft{}The regulations, the rules (of a place)\textquotedblright{}.) Say it once more.
\end{tcolorbox}
\section[obbligatorio]{obbligatorio — compulsory, mandatory}

\label{lesson:IT-C384-obbligatorio}



\begin{quote}
Before the new one: say the Italian for the rule, then the Italian for the regulations.
\end{quote}

\subsection*{You'll want to know: obbligatorio}
\textbf{obbligatorio} — \textquotedblleft{}compulsory, mandatory\textquotedblright{}.

\begin{quote}
In macchina la cintura è obbligatoria.
\end{quote}

\begin{quote}
\emph{In the car the seat belt is compulsory.}
\end{quote}

\emph{Non è obbligatorio} — it is not compulsory.

\begin{cousinweb}
\emph{obbligatorio} is a learned borrowing from Late Latin \emph{obligātōrius} \textquotedblleft{}binding\textquotedblright{}, from \emph{obligāre} \textquotedblleft{}to bind, to oblige\textquotedblright{} (Lewis \& Short), built on \emph{ligāre} \textquotedblleft{}to tie\textquotedblright{}.
\end{cousinweb}

\begin{grammarlens}[title={What you've built}]
One more word for telling, explaining and describing at B1.
\end{grammarlens}

\subsection*{Your turn}
\begin{itemize}
\raggedright
  \item \emph{Say it:} \emph{obbligatorio}
  \item \emph{Say it:} \emph{obbligatorio}, once more
  \item \emph{Say it:} a sentence of your own with \emph{obbligatorio}
  \item \emph{Recall:} say the Italian for the rule, then the Italian for the regulations, then say \emph{obbligatorio} again
\end{itemize}

\begin{tcolorbox}[breakable,colback=teal!4,colframe=teal!35!black,title={Before you move on}]
What does obbligatorio mean? (\textquotedblleft{}Compulsory, mandatory\textquotedblright{}.) Say it once more.
\end{tcolorbox}
\section[facoltativo]{facoltativo — optional}

\label{lesson:IT-C384-facoltativo}



\begin{quote}
Before the new one: say the Italian for the regulations, then the Italian for compulsory.
\end{quote}

\subsection*{You'll want to know: facoltativo}
\textbf{facoltativo} — \textquotedblleft{}optional\textquotedblright{}.

\begin{quote}
La lezione del sabato è facoltativa.
\end{quote}

\begin{quote}
\emph{The Saturday lesson is optional.}
\end{quote}

Built on \emph{facoltà} in the sense of free choice.

\begin{cousinweb}
\emph{facoltativo} is built on \emph{facoltà} \textquotedblleft{}power, option\textquotedblright{}, a learned borrowing from Latin \emph{facultās} \textquotedblleft{}ability, means\textquotedblright{} (Lewis \& Short): left to your own choice.
\end{cousinweb}

\begin{grammarlens}[title={What you've built}]
One more word for telling, explaining and describing at B1.
\end{grammarlens}

\subsection*{Your turn}
\begin{itemize}
\raggedright
  \item \emph{Say it:} \emph{facoltativo}
  \item \emph{Say it:} \emph{facoltativo}, once more
  \item \emph{Say it:} a sentence of your own with \emph{facoltativo}
  \item \emph{Recall:} say the Italian for the regulations, then the Italian for compulsory, then say \emph{facoltativo} again
\end{itemize}

\begin{tcolorbox}[breakable,colback=teal!4,colframe=teal!35!black,title={Before you move on}]
What does facoltativo mean? (\textquotedblleft{}Optional\textquotedblright{}.) Say it once more.
\end{tcolorbox}
\section[l'obbligo]{l'obbligo — the obligation, the duty}

\label{lesson:IT-C384-obbligo}



\begin{quote}
Before the new one: say the Italian for compulsory, then the Italian for optional.
\end{quote}

\subsection*{You'll want to know: l'obbligo}
\textbf{l'obbligo} — \textquotedblleft{}the obligation, the duty\textquotedblright{}.

\begin{quote}
Non hai nessun obbligo di venire.
\end{quote}

\begin{quote}
\emph{You are under no obligation to come.}
\end{quote}

Plural \emph{gli obblighi}. \emph{Avere l'obbligo di} — to be required to.

\begin{cousinweb}
\emph{obbligo} is built on \emph{obbligare} \textquotedblleft{}to oblige\textquotedblright{}, a learned borrowing from Latin \emph{obligāre} \textquotedblleft{}to bind\textquotedblright{} (Lewis \& Short), built on \emph{ligāre} \textquotedblleft{}to tie\textquotedblright{}.
\end{cousinweb}

\begin{grammarlens}[title={What you've built}]
One more word for telling, explaining and describing at B1.
\end{grammarlens}

\subsection*{Your turn}
\begin{itemize}
\raggedright
  \item \emph{Say it:} \emph{l'obbligo}
  \item \emph{Say it:} \emph{l'obbligo}, once more
  \item \emph{Say it:} a sentence of your own with \emph{l'obbligo}
  \item \emph{Recall:} say the Italian for compulsory, then the Italian for optional, then say \emph{l'obbligo} again
\end{itemize}

\begin{tcolorbox}[breakable,colback=teal!4,colframe=teal!35!black,title={Before you move on}]
What does l'obbligo mean? (\textquotedblleft{}The obligation, the duty\textquotedblright{}.) Say it once more.
\end{tcolorbox}
